\documentclass[10pt,a4paper]{amsart}
\usepackage[margin=19mm]{geometry}
\usepackage{amsmath,amssymb,mathtools}
\usepackage[scr=rsfs,cal=euler]{mathalfa}
\usepackage{xfrac}
\usepackage[unicode,hidelinks,bookmarks=false]{hyperref}
\newcommand{\ie}{i.e.\ }
\newcommand{\cf}{cf.\ }
\newcommand{\resp}{respectively\ }
\newcommand{\Subsection}{\subsection}
\providecommand{\nobreakdash}{\nobreak\hbox{-}}
\setlength{\emergencystretch}{2em}
\multlinegap0pt
\usepackage{bm}
\usepackage{bbm}
\usepackage{dsfont}
\usepackage{xspace}
\usepackage{cancel}
\usepackage{mathtools}
\usepackage{cleveref}
\usepackage{amsthm}
\makeatletter
\@ifundefined{thm}{\newtheorem{thm}{Thm}}{}
\@ifundefined{cor}{\newtheorem{cor}[thm]{Corollary}}{}
\@ifundefined{defn}{\newtheorem{defn}[thm]{Definition}}{}
\@ifundefined{fact}{\newtheorem{fact}[thm]{Fact}}{}
\@ifundefined{lem}{\newtheorem{lem}[thm]{Lemma}}{}
\makeatother
\title[On Algorithmic Robustness of Corrupted Markov Chains]{On Algorithmic Robustness of Corrupted Markov Chains}
\author{Jason Gaitonde, Elchanan Mossel}
\date{Source: arXiv:2507.15176v1 (CC BY 4.0)}
\begin{document}
\maketitle

\noindent\emph{Source and licence.} On Algorithmic Robustness of Corrupted Markov Chains. Version arXiv:2507.15176v1 (CC BY 4.0), distributed under the Creative Commons Attribution 4.0 International licence, as recorded in the arXiv licence metadata for this exact version and shown on the abstract page https://arxiv.org/abs/2507.15176v1. Full rights notice: licenses/arxiv-2507.15176-CC-BY-4.0.txt.\par\smallskip

\noindent\textit{Editorial scope.} Counted unit: the Cor whose source label is \texttt{cor:corrupted\_close} in the article (source0.tex lines 435--444, source label \texttt{cor:corrupted\_close}), delivered together with the article's own complete original proof of it (source0.tex lines 445--502, one \texttt{proof} environment of 58 lines). No mathematics is written, repaired or re-derived: statement and proof are the article's own text line for line. Required context retained uncounted: def:corruption (lines 164--171); fact:contract (lines 262--265); lem:coupling\_main (lines 364--370); cor:pr\_close (lines 384--391). Every definition, display and label of those lines is retained unchanged. Omitted from this package and not redistributed: 1--163, 172--261, 266--363, 371--383, 392--434, 503--570. Nothing inside a retained excerpt is omitted or altered: every retained body is delivered exactly as the article's lines are written, so no omission receipt of any kind accompanies this package. Citations: the retained excerpts carry no citation command at all, so no bibliography entry of the article is transcribed and no reference is redistributed. Reference adaptation: none was needed and none was made; every reference inside the delivered unit resolves inside it, so no unresolved reference or citation is produced. Printed slips kept literal: no printed word, symbol, hypothesis or number of the counted statement or of its proof was altered. Layout and class: the article's own class is \texttt{article} with packages including geometry, amssymb, amsmath, amsthm, arydshln, verbatim, titlesec, subfigure, algorithm2e, bm, pdfpages, wrapfig, hyperref, fontenc; the wrapper is the lane's pinned \texttt{amsart} editorial wrapper at 10pt on A4, which restates the theorem-like environments the retained text needs and adds the article's own macro definitions that the retained text needs (none), with extra wrapper packages (bm, bbm, dsfont, xspace, cancel, mathtools, cleveref). The delivered numbering is the unit's own. Assets: no retained span contains a figure environment, a graphics-inclusion command, a PDF-page inclusion, a table or a TikZ picture; no image file is copied into this package, the package declares no asset dependency and no image file is redistributed with it. Licence: version arXiv:2507.15176v1 of this article is distributed under the Creative Commons Attribution 4.0 International licence, as recorded in the arXiv licence metadata for this exact version and shown on the abstract page https://arxiv.org/abs/2507.15176v1. A full rights notice is delivered at licenses/arxiv-2507.15176-CC-BY-4.0.txt. Characters: every retained excerpt is pure ASCII, so the delivered engine, XeLaTeX with its default Latin Modern text font, reports no missing character.
\begin{defn}[Corrupted Markov Chains]
\label{def:corruption}
    Let $\mathcal{X}$ be a finite state space and let $\mathsf{P}$ be a Markov transition matrix on $\mathcal{X}$ with stationary measure $\pi$. A transition matrix $\widetilde{\mathsf{P}}$ is a \textbf{$\varepsilon$-corruption} of $\mathsf{P}$ if it holds that
    \begin{equation*}
        \sum_{x,y\in \mathcal{X}}\pi(x)\vert \mathsf{P}(x,y)-\widetilde{\mathsf{P}}(x,y)\vert\leq \varepsilon.
    \end{equation*}
    We write $\widetilde{\pi}$ for a stationary distribution of the corrupted chain.
\end{defn}

\begin{fact}
\label{fact:contract}
    $\pi$ is the stationary distribution of $\mathsf{P}$, then $\|\mathsf{P}\|_{p\to p}=1$ for all $p\geq 1$.
\end{fact}

\begin{lem}
\label{lem:coupling_main}
    Let $\mathsf{P}$ be any Markov chain with a stationary distribution $\pi$. Then for any distribution $\nu$ and any $t\in \mathbb{N}$ it holds that
    \begin{equation*}
        \|\pi-\nu\|_1\leq \|\pi -\nu\mathsf{P}^t\|_1+t\|\nu-\nu\mathsf{P}\|_1.
    \end{equation*}
\end{lem}

\begin{cor}
\label{cor:pr_close}
Suppose that a Markov chain $\mathsf{P}$ has spectral gap at least $\gamma$ with stationary distribution $\pi$. For any $\delta\in [0,1]$ and distribution $\mu$, consider the PageRank chain $\mathsf{P}(\delta) = (1-\delta)\mathsf{P}+\delta\bm{1}^T\mu$ and let $\pi_{\delta}$ denote the stationary distribution. Then 
\begin{equation*}
    \|\pi-\pi_{\delta}\|_1\leq \frac{C\delta\log\left(\frac{\gamma\|\mu/\pi\|_{\infty}}{\delta}\right)}{\gamma},
\end{equation*}
for an absolute constant $C>0$.
\end{cor}

\begin{cor}
\label{cor:corrupted_close}
    Let $\mathsf{P}$ be a Markov chain with stationary distribution $\pi$ and suppose that $\widetilde{\mathsf{P}}$ is a $\varepsilon$-corruption. For any $\delta>0$ and distribution $\mu$, let $\mathsf{P}(\delta)$ and $\widetilde{\mathsf{P}}(\delta)$ denote the PageRank version with restart $\mu$ with stationary distribution $\pi_{\delta}$ and $\widetilde{\pi}_{\delta}$ as before. Then if $\mu$  is $(\beta,p)$-smooth with respect to $\pi$, 

    \begin{equation*}
        \|\pi_{\delta}-\widetilde{\pi}_{\delta}\|_1\leq \frac{C\varepsilon^{1/q}\log(1/\varepsilon)\beta}{\delta}. 
    \end{equation*}
    for an absolute constant $C>0$, where $q$ is the dual exponent of $p$.

\end{cor}

\begin{proof}
    We again apply \Cref{lem:coupling_main}, but now replacing the role of $\mathsf{P}$ with $\widetilde{\mathsf{P}}_{\delta}$:
    \begin{align*}
        \|\pi_{\delta}-\widetilde{\pi}_{\delta}\|_1&\leq \| \widetilde{\pi_{\delta}}-\pi_{\delta}\widetilde{\mathsf{P}}(\delta)^t\|_1+t\|\pi_{\delta}-\pi_{\delta}\widetilde{\mathsf{P}}(\delta)\|_1.
    \end{align*}
    For the first term, we can use a standard coupling argument as in [Corollary 5.5., Levin-Peres] to assert that 
    \begin{equation*}
        \| \widetilde{\pi_{\delta}}-\pi_{\delta}\widetilde{\mathsf{P}}(\delta)^t\|_1\leq 2\Pr(\tau_{\delta}\geq t)\leq 2\exp(-\delta t),
    \end{equation*}
    where $\tau_{\delta}$ is a geometric random variable with success probability $\delta$ since the total variation distance is bounded by the time that it takes to choose to randomly restart both chains in an optimal coupling starting with $\widetilde{\pi}_{\delta}$ and $\pi_{\delta}$ when transitioning via $\widetilde{\mathsf{P}}$. For the second term, stationarity of $\pi_{\delta}$ with respect to $\mathsf{P}(\delta)$ implies:
    \begin{align*}
        \|\pi_{\delta}-\pi_{\delta}\widetilde{\mathsf{P}}(\delta)\|_1&=\|\pi_{\delta}\mathsf{P}(\delta)-\pi_{\delta}\widetilde{\mathsf{P}}(\delta)\|_1\\
        &=(1-\delta)\|\pi_{\delta}(\mathsf{P}-\widetilde{\mathsf{P}})\|_1\\
        &\leq \sum_{x,y}\pi_{\delta}(x)\vert \mathsf{P}(x,y)-\widetilde{\mathsf{P}}(x,y)\vert.
    \end{align*}
    To conclude, suppose that $\mu$ is $(\beta,p)$-smooth with respect to $\pi$. In this case, using H\"older's inequality gives
    \begin{align*}
        \sum_{x,y}\pi_{\delta}(x)\vert \mathsf{P}(x,y)-\widetilde{\mathsf{P}}(x,y)\vert&=2\mathbb{E}_{X\sim \pi_{\delta}}[d_{\mathsf{TV}}(\mathsf{P}(X,\cdot),\widetilde{\mathsf{P}}(X,\cdot))]\\
        &=2\mathbb{E}_{X\sim \pi}\left[\frac{\mathrm{d}\pi_{\delta}}{\mathrm{d}\pi}(X)d_{\mathsf{TV}}(\mathsf{P}(X,\cdot),\widetilde{\mathsf{P}}(X,\cdot))\right]\\
        &\leq 2\left\|\frac{\pi_{\delta}}{\pi}\right\|_{p,\pi}\|d_{\mathsf{TV}}(\mathsf{P}(X,\cdot),\widetilde{\mathsf{P}}(X,\cdot))\|_{q,\pi}\\
        &\leq 2\varepsilon^{1/q}\left\|\frac{\pi_{\delta}}{\pi}\right\|_{p,\pi}
    \end{align*}
    where in the last step we use the general definition of an $\varepsilon$-corruption with respect to $\pi$ to see that
    \begin{equation*}
        \|d_{\mathsf{TV}}(\mathsf{P}(X,\cdot),\widetilde{\mathsf{P}}(X,\cdot))\|_{q,\pi}=\left(\mathbb{E}_{\pi}[d^q_{\mathsf{TV}}(\mathsf{P}(X,\cdot),\widetilde{\mathsf{P}}(X,\cdot))]\right)^{1/q}\leq \varepsilon^{1/q},
    \end{equation*}
    since the total variation is at most $1$ surely.
    But by the same argument as in \Cref{cor:pr_close} using the contractivity of $\mathsf{P}^*$ in $L^p(\pi)$ by \Cref{fact:contract}, it follows that
    \begin{equation*}
        \left\|\frac{\pi_{\delta}}{\pi}\right\|_{p,\pi}\leq \left\|\frac{\mu}{\pi}\right\|_{p,\pi}\leq \beta,
    \end{equation*}
    by assumption. Putting this together, we find that
    \begin{equation*}
        \|\pi_{\delta}-\widetilde{\pi}_{\delta}\|_1\leq 2\exp(-\delta t)+2\varepsilon^{1/q} \beta t .
    \end{equation*}
    
    We now set $t=\frac{\log(1/\varepsilon)}{\delta}$ to deduce that 
    \begin{equation*}
        \|\pi_{\delta}-\widetilde{\pi}_{\delta}\|_1\leq \frac{C\varepsilon^{1/q}\log(1/\varepsilon)\beta}{\delta}.
    \end{equation*}

\iffalse
    Now suppose instead that $\mu$ is $(\beta,\varepsilon)$-smooth. Recall from \Cref{def:corruption} that
    \begin{equation*}
        \mathbb{E}_{X\sim \pi}[d_{\mathsf{TV}}(\mathsf{P}(X,\cdot),\widetilde{\mathsf{P}}(X,\cdot))]\leq \varepsilon/2.
    \end{equation*}
    Let $\mathcal{S}$ denote the good subset of states $x\in \mathcal{X}$ such that $d_{\mathsf{TV}}(\mathsf{P}(X,\cdot),\widetilde{\mathsf{P}}(X,\cdot))\leq \sqrt{\varepsilon}/2$. Markov's inequality then implies that
    \begin{equation*}
        \pi(\mathcal{S})\leq \sqrt{\varepsilon}.
    \end{equation*}
    By the definition of $(\beta,\varepsilon)$-smoothness, it follows that
    \begin{equation*}
        \mu(\mathcal{S})\leq \beta\sqrt{\varepsilon}
    \end{equation*}
    and so it follows that 

    \fi
\end{proof}

\end{document}
